% !PRINTED 91
\weberpage{91}
%<<CONT

\begin{equation}
 B=\begin{vmatrix}
 a_1^{(1)}&a_2^{(1)}&\cdots&a_\nu^{(1)}\\
 a_1^{(2)}&a_2^{(2)}&\cdots&a_\nu^{(2)}\\
 \vdots&\vdots&&\vdots\\
 a_1^{(\nu)}&a_2^{(\nu)}&\cdots&a_\nu^{(\nu)}
 \end{vmatrix},
\tag{6}
\end{equation}
ihre Unterdeterminanten mit $B_k^{(h)}$, worin $h,k$ nur die Werthe von $1$ bis $\nu$ durchlaufen, und setzen nun voraus, \gesperrt{alle $(\nu+1)$reihigen Unterdeterminanten von $A$ verschwinden}. Wenn wir die ersten $\nu$ Gleichungen des Systems (1) mit $B_k^{(1)},B_k^{(2)},\ldots,B_k^{(\nu)}$ multipliciren und addiren, so folgt [\S.\,22, (2), (6)]:
\begin{equation}
 Bx_k=\sum_{i=1}^{\nu}B_k^{(i)}y_i-
 \sum_{s=\nu+1}^{n}x_s\sum_{i=1}^{\nu}B_k^{(i)}a_s^{(i)}.
\tag{7}
\end{equation}
Hierdurch sind $\nu$ der Unbekannten $x$ durch die übrigen bestimmt, und wir setzen, um zu sehen, inwieweit hierdurch die Gleichungen (1) befriedigt sind, die Ausdrücke in (1) ein. Wir multipliciren hierzu (7) mit $a_k^{(\lambda)}$, wo $\lambda$ die Werthe $1,2,\ldots,n$ durchläuft, und summiren in Bezug auf $k$. So folgt:
\begin{equation}
 B\sum_{k=1}^{\nu}a_k^{(\lambda)}x_k
 =\sum_{i=1}^{\nu}\sum_{k=1}^{\nu}B_k^{(i)}a_k^{(\lambda)}y_i
 -\sum_{s=\nu+1}^{n}x_s\sum_{i=1}^{\nu}\sum_{k=1}^{\nu}B_k^{(i)}a_s^{(i)}a_k^{(\lambda)}.
\tag{8}
\end{equation}
und wenn wir beiderseits
\[
 B\sum_{k=\nu+1}^{n}a_k^{(\lambda)}x_k
\]
addiren:
\begin{equation}
 B\sum_{k=1}^{n}a_k^{(\lambda)}x_k
 =\sum_{i=1}^{\nu}\sum_{k=1}^{\nu}B_k^{(i)}a_k^{(\lambda)}y_i
 +\sum_{s=\nu+1}^{n}x_s\left(Ba_s^{(\lambda)}-
 \sum_{i=1}^{\nu}\sum_{k=1}^{\nu}B_k^{(i)}a_s^{(i)}a_k^{(\lambda)}\right).
\tag{9}
\end{equation}
Nun ist nach \S.\,23, (12) der Coëfficient von $x_s$ auf der rechten Seite eine $(\nu+1)$reihige Unterdeterminante von $A$ und also nach der Voraussetzung gleich Null. Nach (1) soll
\[
 \sum_{k=1}^{n}a_k^{(\lambda)}x_k=y_\lambda
\]
sein, und folglich genügen die Ausdrücke (7) dann und nur dann den Gleichungen (1), wenn
\begin{equation}
 B y_\lambda=\sum_{i=1}^{\nu}\sum_{k=1}^{\nu}B_k^{(i)}a_k^{(\lambda)}y_i
\tag{10}
\end{equation}
ist. Diese Gleichung ist, wenn $\lambda\le\nu$ ist, wegen \S.\,22, (2), (6) immer befriedigt. Ist aber $\lambda=\nu+1,\nu+2,\ldots,n$, so sind $n-\nu$ Bedingungen für die $y$ in (10) enthalten, die, da jede
%CONT>>
